\section{Related Works}

\paragraph{SST with LLM} Recent work shows that simultaneous speech translation (SST) achieves markedly higher quality when large language models (LLMs) serve as the backbone~\citep{ahmad-etal-2024-findings}. \citet{koshkin-etal-2024-transllama} adapt an LLM to SST via finetuning on synthetic translation trajectories, and \citet{ouyang-etal-2025-infinisst} scale this idea with an interleaving architecture for unbounded speech streams. \citet{fu2025efficientadaptivesimultaneousspeech} explore alternative training strategies and data-synthesis pipelines, while \citet{guo2025streamuniachievingstreamingspeech} unify streaming transcription and translation with a truncation mechanism that prunes history based on transcribed inputs. However, these systems still struggle with domain-specific terminology and \method addresses this by integrating a terminology retriever directly into the SST process.

\paragraph{Retrieval-Augmented NMT} Retrieval is an effective remedy for translating rare words, named entities, and domain-specific terminology in NMT. Early approaches retrieve translation pieces from similar training pairs~\citep{zhang-etal-2018-guiding}, mine bilingual terminology with bilingual embeddings~\citep{huck-etal-2019-better}, or fuse dictionary entries into NMT via a Pointer-Disambiguator-Copier~\citep{zhang-etal-2021-point}. kNN-MT~\citep{khandelwal2021nearest} augments NMT with nearest-neighbor lookup over a large datastore, and follow-up work improves its robustness~\citep{jiang-etal-2021-learning} and efficiency~\citep{meng-etal-2022-fast,zhu-etal-2023-ink}. Other approaches extract and attend over entity candidates~\citep{zeng-etal-2023-extract} or combine retrieval and demonstration with LLMs~\citep{li-etal-2025-leveraging}. For speech translation, \citet{li-etal-2024-optimizing-rare} prepend retrieved speech segments as demonstrations, and \citet{wu-etal-2025-locate} use sliding-window retrieval to inject terms into a speech-LLM prompt. All of these target offline translation, whereas \method integrates retrieval directly into incremental streaming translation.

\paragraph{Contextual-Biasing ASR} Recognizing rare words and domain-specific terminology, commonly known as contextual biasing, is also well studied in ASR. Early work splits into two lines: joint optimization of the ASR model with contextual n-gram embeddings~\citep{CLAS,jain20_interspeech}, and shallow fusion of biased words into end-to-end models~\citep{shallow_fusion}, later extended with prefix-trie-based deep biasing~\citep{le21_interspeech} and unified via deep shallow fusion and trie-constrained biasing~\citep{9383560}. \citet{9687898} support runtime addition of context words through an external memory accessed at decoding. To scale biasing lists from hundreds to tens of thousands of entries, subsequent work leverages neural associative memories~\citep{9747726,10023323,wu23e_interspeech,wu-etal-2024-deferred} and retrieval augmentation~\citep{9413756,10095857,jayanthi-etal-2023-retrieve}. 
More recent efforts build on speech foundation models. \citet{10832154} use coarse CTC decoding to filter hotwords into an LLM-based ASR prompt, and \citet{gong24b_interspeech} inject bias words into a SpeechLLM via prompt formats and a biasing-fusion network. BR-ASR~\citep{gong25_interspeech} decouples bias retrieval from ASR decoding through contrastive learning and curriculum training, scaling to 200k entries at 20ms retrieval latency, while GLCLAP~\citep{kong25_interspeech} retrieves bias words directly from speech via global-local contrastive language-audio pre-training, enabling plug-in biasing on frozen models like Whisper. Finally, PAC~\citep{11462669} combines pronunciation-guided context learning with pronunciation-discriminative reinforcement learning. 
Unlike ASR, where the transcript is time-aligned with the speech, SST translations lag behind the input, so the model must therefore decide whether and when to apply each retrieved term.